\section{Bethe--Salpeter 方程}

Gross 第 27 章：两点响应与四点顶点由不可约核 $I$ 与 Bethe--Salpeter（BS）方程联系；推迟极化传播子即密度--密度响应。有效相互作用 $W$ 与穿衣顶点 $\Lambda$ 见下一节。

\subsection{四点顶点}

定义连通四点函数 $\Gamma(1,2;3,4)$（两条入、两条出）。在粒子--空穴通道，
\begin{equation}
  \Gamma
  =I
  +I\,G G\,\Gamma,
\label{eq:BS}
\end{equation}
$I$ 为该通道不可约核（不能用一条 $G G$ 对切开）。密度响应
\begin{equation}
  \chi
  =\chi_0+\chi_0 I\chi
  =\bigl(1-\chi_0 I\bigr)^{-1}\chi_0.
\end{equation}

\subsection{常见截断}

\begin{itemize}
  \item $I=v$：RPA / 含时 Hartree，集体极点即等离激元；
  \item $I=v-v_{\mathrm{ex}}$（交换梯子）：含时 HF，改善自旋通道；
  \item $I=\delta\Sigma/\delta G$：守恒（Baym--Kadanoff）近似，与 $\Phi$ 泛函一致。
\end{itemize}

\subsection{与线性响应}

推迟极化 $\Pi^R$ 即密度对偶极耦合外势的响应核：$\delta n=\Pi^R v_{\mathrm{ext}}$。时间序 $\Pi$ 与 $\Pi^R$ 由 Lehmann 表示相联系（与单粒子 $G$、$G^R$ 同构）。自由气体的 $\Pi_0$ 即 Lindhard 函数，其闭式见第 5 章。

\begin{example}[从 BS 到 RPA]
无顶点修正、$\Sigma=0$ 时 $I=v$、$\chi_0=\Pi_0$，立刻 $\chi=\Pi_0/(1-v\Pi_0)$，回到 \eqref{eq:RPA}。
\end{example}
